\chapter{Build, test and release}
\label{ch:build}

\section{Building}
You need a compiler. \code{build.bat release all} on Windows finds Visual Studio or the Build Tools; \code{./build.sh
release all} on Linux (g++, libx11-dev, libxi-dev) and macOS (the command line tools); CMake is an optional route.
Each build produces \code{Blendity}, \code{blendity\_tests} and \code{blendity\_stress} in \code{dist/<platform>/}. With
\code{blender/lib/<platform>} beside the project (fetched by \code{packaging/fetch\_blender\_libs.sh} at the commits
pinned in \code{packaging/blender-libs.txt}, about 600 MB) the build links Blender's libraries and copies their
runtime files beside the executables; \code{BLENDITY\_NO\_LIBS=1} forces the dependency-free build. The two builds are
the two test suites the user asked for: both must keep building and passing (C-02). Sanitizer builds need the
dependency-free build, as the prebuilt Jolt has no RTTI. The GPU path tracer's shaders are pre-warmed by the build;
the OptiX kernel's PTX is committed so no CUDA toolkit is needed.

\section{The test suites}
\code{blendity\_tests} runs every test (\code{test("what the user sees", []\{ CHECK(...) \})}) in 16 groups
(\code{--list-groups}, \code{--group <name>}; \code{BLENDITY\_TEST\_FILTER} for one by name): modeling, render,
lighting, gpu, assets, animation, physics, editor, livelink, input, sketch, board, sculpt, hci, vr and main. UI tests
drive the editor headless with synthetic mouse and key events; geometry tests call \code{meshops}; every test runs in
a scratch project. \code{tools/run\_suites.py} runs every group of every build given as its own process, each with its
share of the cores, longest first, and exits 1 on a failure or a crash: both Windows builds in about a minute and a
half.

\begin{center}
\begin{tabular}{lrr}
\toprule
Suite at 0.13.0 & Checks & Tests \\
\midrule
Windows, with Blender's libraries & 15,902 & 1,543 \\
Windows, dependency-free & 15,253 & 1,543 \\
Linux (WSL), dependency-free & 15,259 & 1,543 \\
\bottomrule
\end{tabular}
\end{center}

The dependency-free build skips the library-only tests (Embree, OIDN, Manifold, Jolt, FFmpeg, EXR) and says so; the
counts differ by those. \code{blendity\_stress} is the other suite: limits, naive-against-optimised comparisons, problem
and overlap counts, Live Link's cost, the Blender parity rows; \code{--quick} runs in CI in about a minute, and
\code{stress/results/latest.md} is the committed baseline.

\section{The audits}
\code{tools/test\_coverage\_map.py --check} (no command, action, window, button or field without a test),
\code{tools/hci\_consistency\_audit.py --check} (the interface conventions: ``\ldots'' means asks first, Title Case,
tooltips as sentences, no command names in visible text, labels that fit), \code{tools/hci\_klm.py --check} (the KLM
table recounted), \code{tools/hci\_label\_audit.py} (no hard-coded shortcut label), \code{tools/unity\_meta.py --check}
(the Unity bridge's \code{.meta} files complete), \code{tools/make\_books.py --check} (the reference list and index
tables not drifted, every entry cited). The HCI sweep (\code{tests/test\_hci.cpp}) walks 826 screens for cut labels,
narrow panes and contrast.

\section{Continuous integration}
\code{build.yml} on every push and pull request: windows, linux (ubuntu-latest) and macos build everything, run the
unit tests with each failing check as an annotation, run a quick stress pass under a 15-minute alarm, take a headless
screenshot and upload the executables; linux-sanitizers builds the tests with ASan + UBSan and runs them. All four are
required to merge. \code{release.yml} on a \code{v*} tag, on a merge into \code{main} and on a pull request touching
packaging: builds Blendity-full and Blendity on each system with the libraries fetched and cached by their pinned
commits, checks the builds are what they should be (\code{--probe}), installs and uninstalls on Windows (full and lite
types), starts the AppImage and the mounted \code{.dmg} headless, packages the five bridges, builds the two books
(\code{make\_books.py --check --pdf}), and on a tag publishes. \code{moonray.yml} builds MoonRay for the bridge.

\section{Packaging and deployment}
\begin{figure}[H]
  \centering
  \resizebox{\textwidth}{!}{% Packaging and deployment: from a tag to the installed program and its updater.
\begin{tikzpicture}[node distance=5mm and 7mm]
  \node[box] (tag) {A \code{v*} tag\\ \footnotesize pushed when the user asks};
  \node[box, right=of tag, minimum width=46mm] (meta) {\code{release.yml} / \code{meta}\\ \footnotesize the version from \code{VERSION};\\ \footnotesize a merge into \code{main} builds\\ \footnotesize artifacts only};
  \node[box, right=of meta, yshift=18mm] (win) {windows\\ \footnotesize Blendity-full + Blendity (no libs),\\ \footnotesize Inno Setup, install / start / uninstall checked};
  \node[box, right=of meta] (lin) {linux (Ubuntu 22.04)\\ \footnotesize AppImage + .tar.gz,\\ \footnotesize GCC 14's runtime shipped};
  \node[box, right=of meta, yshift=-18mm] (mac) {macos (arm64)\\ \footnotesize a .dmg, started headless};
  \node[box, right=of lin, yshift=-30mm] (books) {books\\ \footnotesize \code{make\_books.py --check --pdf}};
  \node[box, right=of lin, yshift=30mm] (bridges) {bridges\\ \footnotesize the five Live Link packages\\ \footnotesize (artifacts; attached by hand)};
  \node[box, right=22mm of lin, fill=blAmber!15, draw=blAmber] (rel) {\code{publish\_release.sh}\\ \footnotesize the GitHub Release: 3 installers under fixed\\ \footnotesize names, \code{SHA256SUMS.txt}, the two PDFs;\\ \footnotesize the libraries archives into the\\ \footnotesize ``libraries'' release};
  \node[box, right=of rel] (user) {The user's machine\\ \footnotesize \code{releases/latest/download/<name>};\\ \footnotesize Blendity hands over to Blendity-full\\ \footnotesize when the libraries load};
  \node[box, below=of user] (upd) {The updater (\code{src/core/releases.cpp})\\ \footnotesize once a day: a newer release? Install\\ \footnotesize checks the SHA-256; Version History\\ \footnotesize installs any build, older ones too};
  \draw[flow] (tag) -- (meta);
  \draw[flow] (meta) -- (win);
  \draw[flow] (meta) -- (lin);
  \draw[flow] (meta) -- (mac);
  \draw[flow] (meta) |- (books);
  \draw[flow] (meta) |- (bridges);
  \draw[flow] (win) -- (rel);
  \draw[flow] (lin) -- (rel);
  \draw[flow] (mac) -- (rel);
  \draw[flow] (books) -- (rel);
  \draw[flow] (rel) -- (user);
  \draw[flow, <->] (user) -- (upd);
  \draw[dflow] (upd.west) -| node[lbl, pos=0.25] {\code{gh api releases}} (rel.south);
\end{tikzpicture}
}
  \caption{From a tag to the installed program and its updater.}
  \label{fig:deploy}
\end{figure}

Each system gets one installer holding Blendity (dependency-free, always starts) and Blendity-full (with the
libraries): an Inno Setup installer for Windows \cite{InnoSetup2026} installing for the user alone with a \code{.scene}
association, an AppImage \cite{AppImage2026} and a plain \code{.tar.gz} for Linux built on the oldest supported Ubuntu,
a \code{.dmg} for Apple-silicon Macs. \code{publish\_release.sh} publishes the three installers under names without the
version (so \code{releases/latest/download/<name>} always fetches the newest), \code{SHA256SUMS.txt} for everything, the
two PDFs, and the libraries archives into one ``libraries'' release that Blendity's updater reads; a \code{0.x} version
is titled a pre-release of the program, a suffixed tag (\code{-beta.1}) a GitHub pre-release. The release notes are
GitHub's summary of the merged pull requests followed by \code{packaging/install-notes.md}.

The updater (\adr{0033}) asks GitHub once a day, when online and when the user allows, whether a newer release exists,
shows what it brings with Install / Not now / Skip this one, checks the download against the SHA-256 list, saves the
scene and installs; Help $>$ Version History installs any build, older ones too. An install without the libraries
offers to download them for this build; Preferences $>$ Libraries is always the user's choice.

\section{Releasing}
A release is the user's decision. The steps: \code{VERSION} bumped (with the fallback in \code{releases.cpp}, the
installer script's example, the README's pre-release note and the manual's mention), a CHANGELOG entry, a merge,
\code{git tag vX.Y.Z} on the merge and \code{git push origin vX.Y.Z}; \code{release.yml} does the rest and the user
checks the assets on the Releases page. Version 0.13.0 is the release this book was written for; versions before 1.0
are pre-releases.
